\subsubsection*{Resultados del LightGBM}

En este apartado vamos a mostrar los resultados y comentaremos un poco  al respecto, aunque ahondaremos más en la parte de IA explicable (XAI).

\subsubsection*{Modelo baseline}

EVALUACIÓN: LightGBM - Baseline

MÉTRICAS PRINCIPALES:
\begin{itemize}
	\item Accuracy:  1.0000  (Porcentaje total de aciertos)
	\item Precision: 1.0000  (De los que dije que eran reales, ¿cuántos lo eran?)
	\item Recall:    1.0000  (De los reales, ¿cuántos detecté?)
	\item F1-Score:  1.0000  (Balance precision-recall)
	\item AUC-ROC:   1.0000  (Capacidad discriminativa)
	\item MCC:       1.0000  (Coeficiente de correlación de Matthews)
\end{itemize}

CLASSIFICATION REPORT:
              precision    recall  f1-score   support

    Real (1)     1.0000    1.0000    1.0000        10
      IA (0)     1.0000    1.0000    1.0000        10

    accuracy                         1.0000        20
   macro avg     1.0000    1.0000    1.0000        20
weighted avg     1.0000    1.0000    1.0000        20


ANÁLISIS DE ERRORES:
\begin{itemize}
	\item Verdaderos Negativos (Reales bien clasificados): 10
	\item Verdaderos Positivos (IA bien clasificados): 10
	\item Falsos Positivos (Reales clasificados como IA): 0
	\item Falsos Negativos (IA clasificados como Reales): 0
	\item Error tipo I (Falsos Positivos): 0.00%
	\item Error tipo II (Falsos Negativos): 0.00%
\end{itemize}

Podemos ver que, pese a ser un modelo baseline, ya funciona perfectamente. Esto puede deberse a que los audios, cuando se escuchan, se detectan algunos donde no se entiende lo que está hablando el locutor y otros donde pese a entenderse mejor, por lo que hasta un modelo simple como este ya clasifica bien.

\subsubsection*{Validación Cruzada Estratificada}

Resultados Validación Cruzada (5 folds):
\begin{itemize}
	\item Accuracy: 0.9875 (+/- 0.0500)
	\item F1-Score: 0.9882 (+/- 0.0471)
	\item AUC-ROC:  1.0000 (+/- 0.0000)
\end{itemize}
   
Estos resultados arrojan una pequeña desviación en exactitud y en f1, lo cual puede indicar que en el modelo baseline el split con el que se entrenó era ``fácil'' y por eso lo hace tan bien, ya que no en todos los splits lo hace perfecto.

\subsubsection*{Grid Search}

En primer lugar veamos cuáles son los mejores parámetros del grid que hemos definido.

\begin{itemize}
	\item num\_leaves: 15
	\item learning\_rate: 0.01
	\item n\_estimators: 100
	\item min\_child\_samples: 20
	\item subsample: 0.6
	\item colsample\_bytree: 0.6
	\item reg\_alpha: 0
	\item reg\_lambda: 0
\end{itemize}

Vemos que difieren de los que nosotros definimos previamente en el modelo baseline, y tiene sentido pues la máquina, tras entrenar con todas las posibilidades, ha evaluado este conjunto de parámetros como el mejor, maximizando f1 como vimos. Veamos entonces qué resultados obtenemos con estos parámetros, aunque viendo el resultado del modelo anterior, nos hacemos una idea.

 EVALUACIÓN: LightGBM - Optimizado

MÉTRICAS PRINCIPALES:
\begin{itemize}
	\item Accuracy:  1.0000  (Porcentaje total de aciertos)
	\item Precision: 1.0000  (De los que dije que eran reales, ¿cuántos lo eran?)
	\item Recall:    1.0000  (De los reales, ¿cuántos detecté?)
	\item F1-Score:  1.0000  (Balance precision-recall)
	\item AUC-ROC:   1.0000  (Capacidad discriminativa)
	\item MCC:       1.0000  (Coeficiente de correlación de Matthews)
\end{itemize}

CLASSIFICATION REPORT:
              precision    recall  f1-score   support

    Real (1)     1.0000    1.0000    1.0000        10
      IA (0)     1.0000    1.0000    1.0000        10

    accuracy                         1.0000        20
   macro avg     1.0000    1.0000    1.0000        20
weighted avg     1.0000    1.0000    1.0000        20


ANÁLISIS DE ERRORES:
\begin{itemize}
	\item Verdaderos Negativos (Reales bien clasificados): 10
	\item Verdaderos Positivos (IA bien clasificados): 10
	\item Falsos Positivos (Reales clasificados como IA): 0
	\item Falsos Negativos (IA clasificados como Reales): 0
	\item Error tipo I (Falsos Positivos): 0.00%
	\item Error tipo II (Falsos Negativos): 0.00%
\end{itemize}

Como observamos y como la intuición nos decía, los resultados son iguales que los de antes, todo perfecto. Posiblemente la razón sea la misma: una mezcla entre un buen split y pocos datos y/o datos poco representativos.